% Gauss, Werke X/1 (1917), native72–76, printed65–69.
\SourceLeaf{072}\NoSourceHead
\SourceGroup{XI}
\HistoricTitle{[BEWEIS DES REZIPROZITÄTSSATZES FÜR DIE BIQUADRATISCHEN RESTE, DER AUF DIE KREISTEILUNG GEGRÜNDET IST.]}
\HistoricCitation{[Zettel in Ec3, Kapsel 43.]}
\NumberedItem{1}
\par\indent Im Gebiete der complexen ganzen Zahlen, $m$ complexe ungerade Primzahl, $m'$ Adjuncte, $\left[\dfrac km\right]$ diejenige der vier Einheiten $1,\ i,\ -1,\ -i$, welcher
\[
k^{\frac14(mm'-1)}
\]
\noindent nach dem Modulus $m$ congruent ist.
\[
\left(\frac kM\right)=\text{Product aller }\left[\frac km\right],
\]
\noindent indem man für $m$ sämmtliche einzelne Primfactoren [von $M$] substituirt.
\NumberedItem{2}
\noindent $p$ reelle positive Primzahl $\equiv1\ (\mathrm{mod.}\ 4)=mm'$\par
\noindent $q$ reelle positive Primzahl $\equiv1$\hspace{2em}, von $p$ verschieden und $=nn'$\par
\noindent Die reellen Theile von $m,n\equiv1\ (\mathrm{mod.}\ 4)$\par
\noindent $x$ eine unbestimmt bleibende Grösse\par
\noindent $\mathrm X=\dfrac{x^p-1}{x-1}$\par
\noindent $e$ Primitivwurzel für den Modulus $m$, und zwar eine solche dass
\[
e^{\frac14(p-1)}\equiv i\ (\mathrm{mod.}\ m)
\]
\vfill\noindent X$_1$.\hfill9\par\clearpage

\SourceLeaf{073}\ArithmeticHead{66}{L}{ARITHMETIK.\quad NACHLASS.}
\noindent oder
\[
\left(\frac em\right)=i.
\]
\[
\begin{aligned}
x+ix^e-x^{ee}-ix^{e^3}+x^{e^4}+\text{u.\,s.\,w.}-ix^{e^{p-2}}&=T\\
x^q+ix^{qe}-x^{qee}-ix^{qe^3}+x^{qe^4}+\text{u.\,s.\,w.}-ix^{qe^{p-2}}&=U.
\end{aligned}
\]
\par\noindent\hspace*{\parindent}Dies vorausgesetzt hat man
\[
\begin{aligned}
TT&\equiv-m\bigl(x-x^e+x^{ee}-x^{e^3}+\text{u.\,s.\,w.}-x^{e^{p-2}}\bigr)\ (\mathrm{mod.}\ \mathrm X)\\
&\equiv-mV\ (\mathrm{mod.}\ \mathrm X).
\end{aligned}
\]
\noindent Folgt unmittelbar aus Anwendung der Principien (\textit{Disqu. Ar.} art. 345) unter Benützung der in der \textit{Theoria Res. Biquadr.} artt. 18—20 gegebenen Sätze[*)].

Und da bekanntlich
\[
VV\equiv p\ (\mathrm{mod.}\ \mathrm X),
\]
\noindent so wird
\[
T^4\equiv mmp\equiv m^3m'\ (\mathrm{mod.}\ \mathrm X).
\]
\par\noindent\hspace*{\parindent}Ferner ist (\textit{Disqu. Arithm.} art. 53[**)])
\[
T^q\equiv U\ (\mathrm{mod.}\ q).
\]
\par\noindent\hspace*{\parindent}Auch ist leicht zu übersehen, dass
\[
\left(\frac qm\right)U\equiv T\ (\mathrm{mod.}\ \mathrm X).
\]
\par\noindent\hspace*{\parindent}Aus der Verbindung dieser drei Sätze folgt leicht auf ähnliche Art wie \textit{Theorematis fundamentalis demonstrationes et ampliat. novae}, art. 6[***)]
\[
\left(\frac qm\right)(m^3m')^{\frac14(q-1)}\equiv1\ (\mathrm{mod.}\ q),
\]
\noindent folglich auch nach mod. $n$ und mod. $n'$.
\par\vfill\noindent\rule{.25\linewidth}{.4pt}\par\smallskip
\begingroup\footnotesize
\noindent[*) Werke I, S. 422; II, S. 82ff.]\par
\noindent[**) Werke I, S. 43.]\par
\noindent[***) Werke II, S. 58.]\par
\endgroup\clearpage

\SourceLeaf{074}\ArithmeticHead{67}{R}{BEWEIS DES BIQUADRATISCHEN REZIPROZITÄTSGESETZES.}
\par\indent Diese beiden Congruenzen aequivaliren den Gleichungen
\[
\begin{aligned}
\left(\frac qm\right)\left(\frac mn\right)^3\left(\frac{m'}n\right)&=1\\
\left(\frac qm\right)\left(\frac m{n'}\right)^3\left(\frac{m'}{n'}\right)&=1
\end{aligned}
\]
\noindent oder was dasselbe ist
\[
\begin{aligned}
\left(\frac qm\right)\left(\frac{m'}n\right)&=\left(\frac mn\right)\\
\left(\frac qm\right)\left(\frac{m'}{n'}\right)&=\left(\frac m{n'}\right).
\end{aligned}
\]
\noindent Da nun allgemein
\[
\left(\frac{m'}n\right)\left(\frac m{n'}\right)=1,\quad\left(\frac{m'}{n'}\right)\left(\frac mn\right)=1,
\]
\noindent so wird, wenn man obige Gleichungen resp. mit $\left(\dfrac m{n'}\right)$, $\left(\dfrac mn\right)$ multiplicirt
\[
\left(\frac qm\right)=\left(\frac mn\right)\left(\frac m{n'}\right)=\left(\frac mq\right).
\]
\NumberedItem{3}
\par\indent Der zweite Fall ist, wo $q\equiv3\ (\mathrm{mod.}\ 4)$.
\[
\begin{aligned}
T&=x+ix^e-x^{ee}-ix^{e^3}+\ldots\\
T_1&=x-ix^e-x^{ee}+ix^{e^3}+\ldots\\
U&=x^q+ix^{qe}-x^{qee}-ix^{qe^3}+\ldots\\
U_1&=x^q-ix^{qe}-x^{qee}+ix^{qe^3}+\ldots
\end{aligned}
\]
\[
\left(\frac qm\right)U\equiv T\ (\mathrm{mod.}\ \mathrm X)
\]
\[
U_1\equiv T_1\left(\frac qm\right)\ (\mathrm{mod.}\ \mathrm X)
\]
\[
T^q\equiv U_1\ (\mathrm{mod.}\ q)
\]
\[
TT_1\equiv(-1)^{\frac14(p-1)}p\ (\mathrm{mod.}\ \mathrm X)
\]
\[
T^{q+1}\equiv\left(\frac qm\right)\mathpunct{.}(-1)^{\frac14(p-1)}\mathpunct{.}p\ (\mathrm{mod.}\ \mathrm X;\ q)
\]
\[
T^{q+1}\equiv\left(\frac{-q}m\right)p
\]
\par\vfill\noindent\hfill9*\par\clearpage

\SourceLeaf{075}\ArithmeticHead{68}{L}{ARITHMETIK.\quad NACHLASS.}
\[
T^4\equiv mmp
\]
\[
T^{q+1}\equiv m^{\frac12(q+1)}p^{\frac14(q+1)}\equiv\left(\frac{-q}m\right)p.
\]
\begingroup\small
\begin{center}
\begin{tabular}{@{}c@{\quad}c@{\quad}c@{}}
\multicolumn{1}{l}{Also}&\multicolumn{2}{c}{Also}\\[.7ex]
\begin{tabular}{@{}c@{}}
$m^{\frac12(q+1)}p^{\frac14(q-3)}\equiv\left(\dfrac{-q}m\right)$\\[1ex]
Erhoben zur Potenz $\tfrac12(q-1)$\\[1ex]
$m^{\frac14(qq-1)}p^{\frac18(q-3)(q-1)}\equiv\left(\dfrac{-q}m\right)^{\frac12(q-1)}$
\end{tabular}
&\multicolumn{1}{|c|}{\begin{tabular}{c}$q\equiv3\ (\mathrm{mod.}\ 8)$\\[1ex]$\left(\dfrac mq\right)\equiv\left(\dfrac{-q}m\right)$\\[1ex]\strut\end{tabular}}
&\begin{tabular}{c}$q\equiv7\ (\mathrm{mod.}\ 8)$\\[1ex]$\ast$ unten\\[1ex]?\end{tabular}
\end{tabular}
\end{center}\endgroup
\par\indent In diesem Falle wird obige Gleichung
\[
\left(\frac mq\right)\mathpunct{.}p^{\frac12(q-1)}\equiv\left(\frac{-q}m\right)^3\ (\mathrm{mod.}\ q).
\]
\noindent Zugleich aber ist allgemein (für $q\equiv3\ (\mathrm{mod.}\ 4)$)
\[
\left(\frac mq\right)^2\equiv p^{\frac12(q-1)}\ (\mathrm{mod.}\ q),
\]
\noindent woraus
\[
\left(\frac mq\right)^3\equiv\left(\frac{-q}m\right)^3
\]
\noindent also
\[
\left(\frac mq\right)=\left(\frac{-q}m\right).
\]
\noindent(Man wird $-q=k$ setzen, also
\[
\left(\frac mk\right)=\left(\frac km\right)
\]
\noindent allgemein, wenn $k$ Primzahl $\equiv1\ (\mathrm{mod.}\ 4)$, positiv oder negativ.)
\NumberedItem{4}
\par\indent Zur Anwendung auf zweigliedrige complexe Zahlen sind zuvörderst zwei Sätze zu bemerken.

1) $\left(\dfrac kl\right)=1$, wenn $k,l$ zwei reelle Zahlen ohne gemeinsamen Divisor und $l$ ungerade.
\clearpage

\SourceLeaf{076}\ArithmeticHead{69}{R}{BEWEIS DES BIQUADRATISCHEN REZIPROZITÄTSGESETZES.}
\par\indent 2) $\left(\dfrac il\right)=i^{\frac12(l-1)}$, wenn $l\equiv1\ (\mathrm{mod.}\ 4)$, also
\[
{}=(-1)^{\frac14(l-1)}.
\]
\noindent Es seien
\[
a+bi=m,\quad A+Bi=M
\]
\noindent ohne gemeinsamen Theiler; $a,A$ beide $\equiv1$.
\[
\begin{aligned}
\left(\frac{a+bi}{A+Bi}\right)\left(\frac A{A+Bi}\right)&=\left(\frac{aA+bAi+(A+Bi)(a-bi)}{A+Bi}\right)\\
&=\left(\frac{aA+bB}{A+Bi}\right)=\left(\frac{A+Bi}{aA+bB}\right)
\end{aligned}
\]
\noindent oder da
\[
\left(\frac A{A+Bi}\right)=\left(\frac{Bi}A\right)=\left(\frac iA\right)=i^{\frac12(A-1)},
\]
\noindent so wird
\[
\left(\frac mM\right)i^{\frac12(A-1)}=\left(\frac M{aA+bB}\right)
\]
\noindent und ebenso
\[
\left(\frac{M_1}{m_1}\right)i^{\frac12(a-1)}=\left(\frac{m_1}{aA+bB}\right),
\]
\noindent also multiplicirend
\[
\begin{aligned}
\left(\frac mM\right)\left(\frac{M_1}{m_1}\right)i^{\frac12(a+A-2}&=\left(\frac{Mm_1}{aA+Bb}\right)\\
&=\left(\frac{aA+bB+(aB-bA)i}{aA+bB}\right)\\
&=\left(\frac i{aA+Bb}\right)=i^{\frac12(aA+bB-1)}.
\end{aligned}
\]
\par\noindent\hspace*{\parindent}Folglich
\[
\begin{aligned}
\left(\frac mM\right)\left(\frac{M_1}{m_1}\right)&=i^{\frac12(aA+bB-1-A-a+2)}=i^{\frac12((a-1)(A-1)+bB)}\\
&=i^{\frac12bB}=\varepsilon,
\end{aligned}
\]
\noindent also $\pm1$ je nachdem $\tfrac12b,\tfrac12B$ nicht beide ungerade [oder] beide ungerade. Also
\[
\left(\frac mM\right)=\varepsilon\left(\frac Mm\right)\qquad\text{W. Z. B. W.}
\]
\begin{center}\textbf{BEMERKUNG.}\end{center}
\begingroup\small
\par\indent Vergl. die auf diesen Beweis bezüglichen Ausführungen in dem Artikel 23 des \textsc{Bachmann}schen Aufsatzes »Über \textsc{Gauss}' zahlentheoretische Arbeiten«. \textsc{Schering} bemerkt (Göttinger Nachrichten 1879, S. 384), daß \textsc{Gauss} in einem Briefe an \textsc{Eisenstein} (wie dieser an \textsc{Riemann} und \textsc{Riemann} wieder mündlich an \textsc{Schering} mitgeteilt habe) die Grundzüge seines Beweises des biquadratischen Reziprozitätssatzes mit Hilfe der Kreisteilung angegeben haben soll; leider sind die Briefe von \textsc{Gauss} an \textsc{Eisenstein} nicht wieder aufgefunden worden.
\begin{flushright}\textsc{Schlesinger.}\end{flushright}
\ClosingRule\clearpage\endgroup
